% !TEX root =  main.tex
\chapter{Training Approach}
\label{chap:training_approach}
\pagestyle{plain}

Getting from a pile of validated SQL queries to a neural network that estimates cardinalities accurately involves several non-trivial steps: labeling each query with its true row count, encoding the query into a fixed-size tensor the network can digest, and training the network with an appropriate loss function. This chapter covers each of these steps. The active learning framework that decides which queries are worth labeling is presented separately in Chapter~\ref{chap:active_learning}.


\section{Database Labeling}

Every supervised model needs labels. In our setting a label is the true cardinality of a query the number of rows it returns. Obtaining that number means actually running the query: reconstructing a \texttt{SELECT COUNT(*)} statement from the query's tables, joins, and predicates, submitting it to PostgreSQL, and recording the result.

This sounds straightforward, but at scale it is the dominant bottleneck of the entire pipeline. A single three-table join over tens of millions of rows can take seconds. Multiply that by a few thousand queries and the labeling stage alone can consume hours of compute. Some queries never finish at all: they hit pathological join orderings or blow up intermediate result sizes, making them impractical to execute within any reasonable time budget.

The pipeline handles these cases pragmatically. Each query runs under a timeout. If it does not finish, it is retried once with a longer deadline. If it still fails, it receives a fallback cardinality of one, the smallest value that keeps the log-scale normalization well-defined. This fallback introduces a small amount of noise into the training data, but in practice the fraction of queries that reach it is low enough that the model's overall accuracy is not materially affected.

The expense of labeling is the single strongest argument for active learning. If every label costs real compute, then choosing which queries to label is not just a nice optimization it is the difference between a pipeline that finishes overnight and one that runs for days.


\section{Feature Extraction and Bitmap Generation}

The MSCN architecture expects three fixed-size tensors as input one for tables, one for predicates, one for joins but real queries have variable numbers of each. A single-table query has one table entry and zero joins; a three-table query with five filters has three table entries, two joins, and five predicates. The encoding stage bridges this gap through zero-padding and masking.

First, the system scans every query in the pool and builds four sorted vocabularies: one mapping each unique table name to a one-hot index, one for predicate columns, one for comparison operators, and one for join conditions. Sorting alphabetically guarantees that the same schema always produces the same encoding, regardless of the order in which queries were generated.

Each table entry in a query is encoded by concatenating its one-hot vector with a bitmap feature:

\begin{equation}
\mathbf{s}_i = [\text{onehot}(t_i) \| \text{bitmap}(t_i)] \in \mathbb{R}^{|T| + n_{\text{samples}}}
\end{equation}

The bitmap is where the encoding gets interesting. At startup the pipeline samples a fixed set of rows from each table typically one thousand and holds them in memory. For each query, the system checks which sampled rows satisfy the query's filter predicates and sets the corresponding bit to one. The result is a compact binary vector that tells the model how selective a query's filters are on a small but representative slice of the actual data. Without this signal the model would have to infer selectivity purely from the query structure, which is far harder.

Each predicate $(c_j, op_j, v_j)$ is encoded as:

\begin{equation}
\mathbf{p}_j = [\text{onehot}(c_j) \| \text{onehot}(op_j) \| \text{norm}(v_j)] \in \mathbb{R}^{|C| + |O| + 1}
\end{equation}

where the predicate value is min-max normalized so that all values fall between zero and one. Join conditions are simply one-hot encoded: $\mathbf{j}_k = \text{onehot}(j_k) \in \mathbb{R}^{|J|}$.

To handle the variable sizes, every query's table, predicate, and join tensors are padded to the maximum observed length and paired with a binary mask:

\begin{equation}
\text{mask}(i) = \begin{cases} 1 & \text{if position } i \text{ is a real entry} \\ 0 & \text{if position } i \text{ is padding} \end{cases}
\end{equation}

The masks ensure that the downstream pooling layers average only over real entries, not over the zeros that were inserted to make the tensors rectangular.


\section{Multi-Set Convolutional Network}

MSCN, introduced by Kipf et al.~\cite{DBLP:conf/cidr/KipfKRLBK19}, processes each of the three query components through its own two-layer neural module, pools the results, and combines them for the final prediction. Figure~\ref{fig:mscn_architecture} shows the full architecture.

\begin{figure}[htbp]
\centering
\resizebox{\textwidth}{!}{%
\begin{tikzpicture}[
    node distance=2cm and 3cm,
    block/.style={draw, rectangle, rounded corners=3pt, minimum width=3cm, minimum height=1cm, align=center, font=\small},
    arrow/.style={->, thick, >=stealth},
]
    % Input layer
    \node [block, fill=blue!20] (s_in) {Table Features\\$\mathbf{S} \in \mathbb{R}^{b \times n_t \times d_s}$};
    \node [block, fill=green!20, right=of s_in] (p_in) {Predicate Features\\$\mathbf{P} \in \mathbb{R}^{b \times n_p \times d_p}$};
    \node [block, fill=yellow!20, right=of p_in] (j_in) {Join Features\\$\mathbf{J} \in \mathbb{R}^{b \times n_j \times d_j}$};

    % MLP layers
    \node [block, below=of s_in] (s_mlp1) {Neural Layer 1\\+ Activation + Dropout};
    \node [block, below=of p_in] (p_mlp1) {Neural Layer 1\\+ Activation + Dropout};
    \node [block, below=of j_in] (j_mlp1) {Neural Layer 1\\+ Activation + Dropout};

    \node [block, below=1.5cm of s_mlp1] (s_mlp2) {Neural Layer 2\\+ Activation + Dropout};
    \node [block, below=1.5cm of p_mlp1] (p_mlp2) {Neural Layer 2\\+ Activation + Dropout};
    \node [block, below=1.5cm of j_mlp1] (j_mlp2) {Neural Layer 2\\+ Activation + Dropout};

    % Pooling
    \node [block, fill=gray!15, below=1.5cm of s_mlp2] (s_pool) {Average Pooling};
    \node [block, fill=gray!15, below=1.5cm of p_mlp2] (p_pool) {Average Pooling};
    \node [block, fill=gray!15, below=1.5cm of j_mlp2] (j_pool) {Average Pooling};

    % Concatenation
    \node [block, fill=red!20, below=1.5cm of p_pool] (concat) {Combine All\\$\mathbb{R}^{3h}$};

    % Output
    \node [block, below=1.5cm of concat] (out_mlp) {Output Layer + Activation + Dropout};
    \node [block, fill=orange!20, below=1.5cm of out_mlp] (sigmoid) {Final Prediction\\$\hat{y} \in [0, 1]$};

    % Arrows
    \draw [arrow] (s_in) -- (s_mlp1);
    \draw [arrow] (p_in) -- (p_mlp1);
    \draw [arrow] (j_in) -- (j_mlp1);
    \draw [arrow] (s_mlp1) -- (s_mlp2);
    \draw [arrow] (p_mlp1) -- (p_mlp2);
    \draw [arrow] (j_mlp1) -- (j_mlp2);
    \draw [arrow] (s_mlp2) -- (s_pool);
    \draw [arrow] (p_mlp2) -- (p_pool);
    \draw [arrow] (j_mlp2) -- (j_pool);
    \draw [arrow] (s_pool) -- (concat);
    \draw [arrow] (p_pool) -- (concat);
    \draw [arrow] (j_pool) -- (concat);
    \draw [arrow] (concat) -- (out_mlp);
    \draw [arrow] (out_mlp) -- (sigmoid);

\end{tikzpicture}%
}
\caption{MSCN architecture. Three independent neural modules process tables, predicates, and joins. Their pooled representations are concatenated and mapped to a cardinality estimate in $[0, 1]$.}
\label{fig:mscn_architecture}
\end{figure}

The design exploits a fundamental property of SQL: the meaning of a query does not change if the tables in the FROM clause or the predicates in the WHERE clause are reordered. Processing each set through a shared-weight module followed by mean pooling makes the network invariant to such permutations, which is exactly the right inductive bias for this problem.

Within each module, the first layer projects the input to a hidden dimension $h$ and the second refines it:

\begin{equation}
\mathbf{h}_i^{(1)} = \text{Dropout}\left(\text{ReLU}\left(\mathbf{W}_1 \mathbf{x}_i + \mathbf{b}_1\right), p\right)
\end{equation}

\begin{equation}
\mathbf{h}_i^{(2)} = \text{Dropout}\left(\text{ReLU}\left(\mathbf{W}_2 \mathbf{h}_i^{(1)} + \mathbf{b}_2\right), p\right)
\end{equation}

The masked mean pool then averages only over the real entries, ignoring padding:

\begin{equation}
\mathbf{h}_{\text{pool}} = \frac{\sum_{i=1}^{n} m_i \cdot \mathbf{h}_i^{(2)}}{\sum_{i=1}^{n} m_i}
\end{equation}

The three pooled vectors are concatenated into a single $3h$-dimensional representation and passed through an output network with sigmoid activation:

\begin{equation}
\hat{y} = \sigma\left(\mathbf{W}_{\text{out2}} \cdot \text{Dropout}\left(\text{ReLU}\left(\mathbf{W}_{\text{out1}} [\mathbf{h}_{\text{sample}} \| \mathbf{h}_{\text{pred}} \| \mathbf{h}_{\text{join}}] + \mathbf{b}_{\text{out1}}\right)\right) + \mathbf{b}_{\text{out2}}\right)
\end{equation}

The sigmoid squashes the output to $[0, 1]$, representing the cardinality as a fraction of the observed range.


\section{Normalization and Loss}

True cardinalities range from one to tens of millions. Training on raw counts would be disastrous: a single query returning ten million rows would dominate every gradient update, and the model would never learn to distinguish between a query returning fifty rows and one returning five hundred. Log-scale min-max normalization compresses this range into $[0, 1]$:

\begin{equation}
y_{\text{norm}} = \frac{\log(y) - \log_{\min}}{\log_{\max} - \log_{\min}}
\end{equation}

At inference time the inverse transformation recovers the actual predicted cardinality:

\begin{equation}
\hat{y}_{\text{actual}} = \exp\left(\hat{y}_{\text{norm}} \cdot (\log_{\max} - \log_{\min}) + \log_{\min}\right)
\end{equation}

We train and evaluate with Q-error, the standard metric in cardinality estimation:

\begin{equation}
Q(p, a) = \max\left(\frac{p}{a}, \frac{a}{p}\right)
\end{equation}

A Q-error of one means a perfect estimate. A Q-error of ten means the prediction is off by a factor of ten, whether it overestimates or underestimates. This symmetry is important: a query optimizer cares about the magnitude of the error, not its direction. Q-error is also scale-invariant being wrong by a factor of two carries the same penalty whether the true cardinality is a hundred or a million.
